\documentclass[../main.tex]{subfiles}

\begin{document}

\chapter{Glossaire des termes techniques}
\label{ann:glossaire}

\begin{description}[leftmargin=4cm, style=nextline]
  \item[Cédante] Compagnie d'assurance qui transfère une partie de ses risques à un
        réassureur en contrepartie d'une prime de réassurance.
  \item[Densité d'assurance] Primes d'assurance par habitant, exprimées en USD.
  \item[IBNR] Incurred But Not Reported : provision pour sinistres survenus mais non
        encore déclarés à la compagnie d'assurance.
  \item[Kaopen] Indice d'ouverture du compte de capital \citep{ChinnIto2006}.
  \item[LOCO] Leave-One-Country-Out : procédure de validation croisée adaptée aux
        panels de petite taille, où chaque pays est successivement exclu de
        l'entraînement pour évaluer la généralisation du modèle.
  \item[Ratio S/P] Sinistres sur Primes ; un ratio structurellement inférieur à
        85\,\% indique un résultat technique attractif pour le réassureur.
  \item[Réassurance] Mécanisme par lequel une compagnie d'assurance transfère à un
        réassureur une partie des risques souscrits en contrepartie d'une prime.
  \item[Rétrocession] Transfert par un réassureur d'une partie des risques acceptés
        à un autre réassureur, afin de gérer sa propre exposition.
  \item[SCAR] Scoring par Corrélation Ajustée et Robustesse : modèle de scoring
        multicritère développé dans ce projet pour classer les marchés africains.
  \item[Sprint] Itération de développement de durée fixe dans une approche Agile,
        à l'issue de laquelle un livrable partiel est produit et validé.
  \item[Taux de pénétration] Ratio primes d'assurance / PIB, exprimé en pourcentage.
  \item[ULR] Ultimate Loss Ratio (ratio sinistres à terme) : ratio intégrant le
        développement des réserves sinistres sur plusieurs exercices.
  \item[Zone CIMA] Périmètre des 14 pays membres de la Conférence Interafricaine des
        Marchés d'Assurances, disposant d'un code des assurances unifié.
\end{description}


\end{document}
